One airframe, four motors, four very different flights. This tutorial is about
motor management: the two local file formats, how databases merge, and the
cleanest experiment QtRocket offers --- swap the motor, change nothing else,
and watch total impulse turn into altitude.

\begin{note}[Setup]
An empty directory with both bundled motor files:
\begin{minted}{text}
cp <qtrocket>/tests/data/motors/estes_small.qmd .
cp <qtrocket>/data/Aerotech.rse .
qtrocket-cli tutorial2.cli      # or type along interactively
\end{minted}
\end{note}

\subsection*{Two files, one database}

\begin{console}
\muted{# QtRocket CLI - type 'help' for commands, 'quit' to exit.}
\uin{loaddb estes_small.qmd}
\ok loaddb: 13 new motors from estes_small.qmd (13 total)
\uin{loadmotors Aerotech.rse}
\ok loadmotors: 249 motors from Aerotech.rse
\uin{listmotors D1}
\ok listmotors: 4 shown (filter="D1")
D11  avg=9.4N  Itot=17.49Ns
D12  avg=10.21N  Itot=16.84Ns
D13W  avg=11.284N  Itot=19.183Ns
D15T  avg=14.012N  Itot=19.617Ns
\uin{listmotors E30}
\ok listmotors: 1 shown (filter="E30")
E30T  avg=32.387N  Itot=39.512Ns
\end{console}


\repl{loaddb} reads QtRocket's native \code{.qmd}; \repl{loadmotors} imports
vendor formats --- RockSim \code{.rse} here, RASP \code{.eng} works the same
way. Both \emph{merge} into one session database: the 13 Estes-class motors
and the 249 AeroTech motors are now one searchable pool, which is why the
\code{D1} filter catches entries from both files at once. Filters are plain
case-sensitive substring matches --- \code{E30} pins down one motor family.

\subsection*{A mid-size airframe}

\begin{console}
\uin{newdesign NoseCone name=nose baseRadius=0.013 length=0.09 density=900}
\ok newdesign: root NoseCone id=2
\uin{addpart nose BodyTube name=body innerRadius=0.0121 outerRadius=0.013 length=0.30 density=1000}
\ok addpart: BodyTube id=3 under 2
\uin{addpart body FinSet name=fins finCount=3 rootChord=0.05 tipChord=0.025 span=0.03 sweep=0.02 thickness=0.003 bodyRadius=0.013 density=900 seat=OnSurface parentStation=0 childStation=0}
\ok addpart: FinSet id=4 under 3
\uin{checkdesign}
\ok checkdesign: 3 parts, contiguous, no overlaps
\uin{listparts}
\ok listparts:
  [2] NoseCone "nose"  m=0.0143351 kg
    [3] BodyTube "body"  m=0.0212906 kg
      [4] FinSet "fins"  m=0.0091125 kg
  -- composite: mass=0.0447382 kg, cg_z=-0.209509 m (t=0)
\end{console}


Nothing new grammatically --- a 26\,mm-diameter, 30\,cm tube under a stubby
nose --- but note the composite mass in \repl{listparts}: about
$45$\,g of airframe. Keep that number in mind as the motors get heavier.

\subsection*{The ladder: C11, D12, E30T, F32T}

\begin{console}
\uin{setdrag 0.75}
\ok setdrag: 0.75
\uin{setarea 0.000531}
\ok setarea: 0.000531 m^2
\uin{setmotor C11}
\ok setmotor: C11
\uin{launch}
\ok launch: 1501 steps, t_final=15.000s
  apogee    = 257.089 m @ t=6.600s
  max_speed = 100.609 m/s @ t=0.770s
  downrange = 0.000 m
  landing   = t=15.000s, pos=(0.000, 0.000, -0.419)
CSV /tmp/flight1/qtrocket_run.csv
\uin{save flight_C11.csv}
\ok save: /tmp/flight1/flight_C11.csv
\uin{setmotor D12}
\ok setmotor: D12
\uin{launch}
\ok launch: 2037 steps, t_final=20.360s
  apogee    = 447.674 m @ t=8.140s
  max_speed = 145.186 m/s @ t=1.590s
  downrange = 0.000 m
  landing   = t=20.360s, pos=(0.000, 0.000, -0.350)
CSV /tmp/flight1/qtrocket_run.csv
\uin{save flight_D12.csv}
\ok save: /tmp/flight1/flight_D12.csv
\uin{setmotor E30T}
\ok setmotor: E30T
\uin{launch}
\ok launch: 2603 steps, t_final=26.020s
  apogee    = 721.934 m @ t=8.540s
  max_speed = 308.837 m/s @ t=0.890s
  downrange = 0.000 m
  landing   = t=26.020s, pos=(0.000, 0.000, -0.175)
CSV /tmp/flight1/qtrocket_run.csv
\uin{save flight_E30T.csv}
\ok save: /tmp/flight1/flight_E30T.csv
\uin{setmotor F32T}
\ok setmotor: F32T
\uin{launch}
\ok launch: 3007 steps, t_final=30.060s
  apogee    = 940.544 m @ t=9.770s
  max_speed = 326.873 m/s @ t=1.300s
  downrange = 0.000 m
  landing   = t=30.060s, pos=(0.000, 0.000, -0.255)
CSV /tmp/flight1/qtrocket_run.csv
\uin{save flight_F32T.csv}
\ok save: /tmp/flight1/flight_F32T.csv
\end{console}


Each iteration is the same three commands: \repl{setmotor}, \repl{launch},
\repl{save} (which copies the trajectory CSV before the next flight overwrites
it). The geometry never changes; only the motor does. Four flights,
four apogees:

\begin{center}\small
\begin{tabular}{@{}lrrr@{}}
\toprule
\textbf{Motor} & \textbf{Class} & \textbf{Apogee} & \textbf{Time to apogee}\\
\midrule
C11  & C & $257.1\m$ & $6.6$\,s\\
D12  & D & $447.7\m$ & $8.1$\,s\\
E30T & E & $721.9\m$ & $8.5$\,s\\
F32T & F & $940.5\m$ & $9.8$\,s\\
\bottomrule
\end{tabular}
\end{center}

Each impulse class doubles the total impulse of the one before it --- and the
apogees march accordingly, sub-doubling because drag grows with speed.
\Cref{fig:motor-ladder} overlays the four saved trajectories.

\begin{figure}[htbp]\centering
\begin{tikzpicture}
\begin{axis}[qtplot,
  xlabel={time (s)}, ylabel={altitude (m)},
  xmin=0, ymin=0,
  legend pos=north east]
\addplot table {data/tut2-C11.dat};  \addlegendentry{C11}
\addplot table {data/tut2-D12.dat};  \addlegendentry{D12}
\addplot table {data/tut2-E30T.dat}; \addlegendentry{E30T}
\addplot table {data/tut2-F32T.dat}; \addlegendentry{F32T}
\end{axis}
\end{tikzpicture}
\caption{The motor ladder: one airframe flown on four motors, plotted from the
four \repl{save}d trajectory files. Impulse class doubles rung to rung; drag
keeps the apogees from doubling with it.}
\label{fig:motor-ladder}
\end{figure}

\subsection*{Keep the collection}

\begin{console}
\uin{savedb my_flight_box.qmd}
\ok savedb: 262 motors to my_flight_box.qmd
\uin{savedesign motor_lab.qrd}
\ok savedesign: motor_lab.qrd
\uin{quit}
\ok bye
\end{console}


\repl{savedb} writes the whole merged database --- both files' motors plus
anything an online search added --- to one \code{.qmd}, so the next session
starts with a single \repl{loaddb}. The saved design
(\cref{fig:motor-lab}) records the \emph{last} selected motor by name.

\screenshot{ug-motor-lab.png}{\code{motor\_lab.qrd}: the 26\,mm test airframe.
Same silhouette --- only the motor changed.}{fig:motor-lab}{0.78}

\begin{tip}[What you learned]
\code{.qmd} vs \code{.rse}/\code{.eng}; databases merge rather than replace;
\repl{listmotors} filters are substrings; \repl{setmotor} swaps propulsion
without touching geometry; \repl{save} preserves a flight's CSV before the
next \repl{launch} overwrites it; \repl{savedb} persists the merged motor pool.
\end{tip}
